\section{Introduzione}
\label{sec:intro}

Fra l'aprile del 2009 e il gennaio del 2011 l'autore ha usato \emph{Derive
for Windows}, il sistema di calcolo simbolico della Texas Instruments, come
un taccuino: si scrive un'espressione, si chiede al programma di
semplificarla, risolverla, derivarla o farne il limite, si cambia un
parametro e si guarda che cosa succede. Ne sono rimasti ventitr\'e file
\texttt{.dfw} con $792$ espressioni. \`E matematica sperimentale nel senso
di~\cite{borwein2008}: raramente una dimostrazione, spesso una regolarit\`a
intravista, a volte un errore.

Derive \`e stato ritirato dal commercio nel 2007 e non \`e pi\`u mantenuto; i
suoi file non si aprono con i programmi di calcolo attuali, ma conservano le
espressioni come testo. Il recupero, che
consiste nel decodificare il formato, tradurre le espressioni in notazione
matematica e verificarle numericamente, \`e descritto
nell'appendice~\ref{app:recupero}. Il risultato completo \`e un'edizione
integrale di oltre duecento pagine, in cui ogni espressione \`e riprodotta e
commentata.

Questa nota \`e un'altra cosa: una scelta per temi, pensata per un lettore
matematico. Ogni sezione dice che cosa fa il quaderno, enuncia in forma
pulita quello che c'\`e di vero, lo dimostra o ne abbozza la dimostrazione,
lo colloca nella letteratura e dice che cosa resta aperto. Nessun risultato
\`e presentato come nuovo. L'unica dimostrazione che non abbiamo trovato in
letteratura, quella della sezione~\ref{sec:fattoriali}, \`e elementare.

In sintesi:
\begin{itemize}[leftmargin=1.4em,itemsep=1pt]
  \item la sezione~\ref{sec:semiangolo} ricostruisce l'idea che attraversa
        cinque quaderni, scrivere l'esponenziale con la tangente del
        semiangolo, e la segue fino a Euler, Lambert, Pad\'e e alla
        trasformazione bilineare dei filtri numerici;
  \item la sezione~\ref{sec:fattoriali} spiega una regolarit\`a fra
        fattoriali che dura quattro passi e dimostra che non si ripete;
  \item le sezioni~\ref{sec:media}, \ref{sec:cosinc}, \ref{sec:gradini}
        e~\ref{sec:e} riconducono altri quattro quaderni a fatti noti: il
        teorema della media per l'equazione di Helmholtz, la trasformata di
        Hilbert, la dequantizzazione di Maslov, l'eliminazione dei termini
        dispari dell'errore;
  \item la sezione~\ref{sec:delta} discute un calcolo con un'interpretazione
        che non regge, e lascia aperta la domanda;
  \item la sezione~\ref{sec:mezzo} raccoglie la formula preferita dei
        quaderni, le sue varianti e la sua forma generale;
  \item la sezione~\ref{sec:breve} riassume gli altri quaderni, la
        sezione~\ref{sec:refusi} gli errori.
\end{itemize}

\paragraph{Come si leggono i riferimenti.}
La scrittura \riga{factor.dfw}{4} indica l'espressione numero~4 del file
\texttt{factor.dfw}, con la numerazione di Derive; \righe{gauss.dfw}{1}{5}
indica un intervallo di espressioni. L'edizione integrale le riproduce tutte
con la stessa numerazione.

\paragraph{Verifica.}
Ogni valore numerico citato nella nota \`e ricalcolato da uno script incluso
nel repository, \texttt{verifica\_nota.py}, con \texttt{mpmath} a sessanta
cifre; un secondo script controlla che ogni decimale presente nel testo sia
fra quelli verificati, e che la versione italiana e quella inglese riportino
gli stessi numeri, gli stessi riferimenti e le stesse citazioni.
